\vspace{-4ex}

\subsubsection{Aggregation}

\paragraph{Conditional aggregation (pivot).} The signal is a request for one
column per category on a single row. Fold a categorical column into counts by
wrapping a \texttt{CASE} in an aggregate. Rows the \texttt{CASE} does not match
contribute 0 to \texttt{SUM}. In Postgres the \texttt{FILTER} clause does the
same thing and reads more cleanly.

\begin{lstlisting}[style=sql, caption={Pivot risk categories to columns}]
SELECT
    inspection_type,
    COUNT(*) FILTER (WHERE risk_category = 'Low Risk')      AS low_risk,
    COUNT(*) FILTER (WHERE risk_category = 'High Risk')     AS high_risk,
    COUNT(*)                                                AS total
FROM sf_restaurant_health_violations
GROUP BY inspection_type
ORDER BY total DESC;
\end{lstlisting}

\vspace{-1ex}

The \texttt{SUM(CASE WHEN ... THEN 1 ELSE 0 END)} form is equivalent and
portable across engines. \texttt{COUNT} ignores \texttt{NULL}, so
\texttt{COUNT(CASE WHEN cond THEN 1 END)} (no \texttt{ELSE}) also counts only
the matches.
